\documentclass[12pt]{article}
\usepackage[utf8]{inputenc}
\usepackage[T1]{fontenc}
\usepackage{hyperref}
\title{Kanakanabu Language Resources for K-12 Curriculum Development}
\author{Compiled by Claude Sonnet 5.5 (claude-sonnet-5-5) for Andr\'as Kornai}
\date{}
\begin{document}
\maketitle

\section{Introduction}
Kanakanabu, also spelled Kanakanavu (ISO 639-3 xnb), is a Southern Tsouic Austronesian language of Taiwan. It is critically endangered. Speaker counts range from 4 (2012) to 6 fluent speakers over age 60 (2014). The Kanakanavu live in Manga and Takanua, Namasia District, Kaohsiung. The group and language were officially recognised in 2014. The search found few resources, and the exact URL of the online dictionary was not found.

\section{Primary Educational Resources}
I found no school textbooks or primers in the search results. The Council of Indigenous Peoples set a Kanakanavu orthography in 2005, and an online dictionary using it is reported. I did not confirm its address.

\section{Linguistic and Reference Materials}
\begin{itemize}
\item Wikipedia, Kanakanavu language: \url{https://en.wikipedia.org/wiki/Kanakanavu_language}
\item Wikipedia, Kanakanavu people: \url{https://en.wikipedia.org/wiki/Kanakanavu_people}
\item Omniglot, Kanakanavu: \url{https://omniglot.com/writing/kanakanavu.htm}
\item Academia Sinica Institute of Linguistics publication, PDF 1: \url{https://www.ling.sinica.edu.tw/upload/researcher_manager_result/ec1a38be221f614d10ce7ba7149fe26d.pdf}
\item Academia Sinica Institute of Linguistics publication, PDF 2: \url{https://www.ling.sinica.edu.tw/upload/researcher_manager_result/c95d71f0cb1814f6db3feee484ef7dce.pdf}
\item The Language Closet, Languages of Taiwan: Kanakanabu: \url{https://thelanguagecloset.com/2021/03/13/languages-of-taiwan-kanakanabu-kanakanavu/}
\end{itemize}

\section{Religious and Community Resources}
No community organisations or scripture translations turned up in the search.

\section{Digital and Audio Resources}
\begin{itemize}
\item Verbix endangered languages page: \url{https://docs.verbix.com/EndangeredLanguages/Kanakanavu}
\end{itemize}

\section{Curriculum Development Context}
There are very few speakers and no known school teaching. A curriculum would depend on the 2005 orthography, the Academia Sinica linguistic work and the online dictionary, once located. Work with the Kanakanavu community in Namasia would be essential. The contents of the Academia Sinica PDFs were not examined.
\end{document}